Un loueur de bicyclettes propose deux types de bicyclettes : des bicyclettes traditionnelles et des bicyclettes électriques.

Il incite ses clients à prendre une assurance.

On dispose des informations suivantes.
\begin{itemize}[label=$\bullet$]
	\item 60\,\% des clients ont loué une bicyclette traditionnelle, les autres ont loué une bicyclette électrique.
	\item Parmi ceux qui ont loué une bicyclette traditionnelle, 25\,\% ont pris une assurance.
	\item 20\,\% de l'ensemble des clients ont pris une assurance.
\end{itemize}

On choisit un client au hasard et on note les événements :

$T$ : \og{}le client a loué une bicyclette traditionnelle\fg{}.

$A$ : \og{}le client a pris une assurance\fg{}.

\begin{enumerate}
	\item Recopier l'arbre pondéré ci-dessous et compléter les pointillés.

	\begin{center}
	\begin{tikzpicture}[xscale=1,yscale=1,rotate=0,every node/.style={scale=1,rotate=0}]
		\node(n0)at(0,0){};
		\node(n1_1)at(3,1){$T$};
		\node(n1_2)at(3,-1){$\overline{T}$};
		\node(n2_1)at(6,1.5){$A$};
		\node(n2_2)at(6,0.5){$\overline{A}$};
		\node(n2_3)at(6,-0.5){$A$};
		\node(n2_4)at(6,-1.5){$\overline{A}$};
		\draw(n0.east)--(n1_1)node[pos=0.6,fill=white]{\dots};
		\draw(n0.east)--(n1_2)node[pos=0.4,fill=white]{\dots};
		\draw(n1_1.east)--(n2_1)node[pos=0.6,fill=white]{\dots};
		\draw(n1_1.east)--(n2_2)node[pos=0.4,fill=white]{\dots};
		\draw(n1_2.east)--(n2_3);
		\draw(n1_2.east)--(n2_4);
	\end{tikzpicture}
	\end{center}

	\item Donner, par simple lecture de l'énoncé, la probabilité de l'événement $A$.

	\item Montrer que la probabilité que le client ait loué une bicyclette traditionnelle et qu'il ait pris une assurance est égale à \np{0.15}.

	\item En déduire que la probabilité \mbox{$P\big(\,\overline{T}\cap A\,\big)$} est égale à \np{0.05}.

	\item Déterminer la probabilité que le client ait pris une assurance sachant qu'il a loué une bicyclette électrique.

	On donnera le résultat sous la forme d'une fraction irréductible.
\end{enumerate}
